% GENERATED FILE. Edit canonical lessons, then run npm run generate:books.
% canonical-source-hash: fnv1a64:c8092733f9f524c6
% canonical-lessons: ZH-C165-dayi, ZH-W165-yi-clothes, ZH-C165-maoyi, ZH-C165-tiao, ZH-C165-qunzi, ZH-C165-weijin

\chapter{Coats and Scarves}
\label{ch:zh-coats-and-scarves}
\hlchaptermodality{165}

\begin{chapteropening}
\textbf{By the end of this chapter:} \emph{I can name a coat, a sweater, a skirt and a scarf, and count them.}

Complete the last lesson of chapter 165: I can name a coat, a sweater, a skirt and a scarf, and count them.
\end{chapteropening}

\section[\texorpdfstring{\zh{大衣} (dàyī)}{大衣 (dàyī)}]{\zh{大衣} (dàyī) — overcoat; coat}

\label{lesson:ZH-C165-dayi}



\begin{quote}
Before the new one: say the Mandarin for to put on make-up, then the Mandarin for to have a shave.
\end{quote}

\subsection*{You'll want to know: \zh{大衣}}
\textbf{\zh{大衣}} — \emph{dàyī} — \textquotedblleft{}overcoat; coat\textquotedblright{}.

\begin{quote}
\textbf{\zh{我的大衣在哪儿}} \emph{wǒ de dàyī zài nǎr} — Where's my coat?
\end{quote}

The new character, in pieces: \zh{衣} — six strokes: a dot and a horizontal on top, then a long stroke falling to the left, a short vertical that flicks up to the right, a short stroke falling to the left, and a long stroke falling to the right. The next lesson puts it in your hand.

\begin{cousinweb}
\textbf{\zh{大衣}} (\emph{dàyī}) is built from \zh{大} 'big' and \zh{衣} 'garment', both inherited from Middle Chinese, \emph{daH} and \emph{'j+j} (Guangyun, in Baxter's transcription), which glosses the second as the garment of the upper body: the big garment over the rest.
\end{cousinweb}

\begin{grammarlens}[title={What you've built}]
One more word for everyday things.
\end{grammarlens}

\subsection*{Your turn}
\begin{itemize}
\raggedright
  \item \emph{Say it:} \emph{dàyī}
  \item \emph{Say it:} \emph{dàyī}, once more
  \item \emph{Say it:} \emph{wǒ de dàyī zài nǎr} — Where's my coat?
  \item \emph{Recall:} say the Mandarin for to put on make-up, then the Mandarin for to have a shave, then say \emph{dàyī} again
\end{itemize}

\begin{tcolorbox}[breakable,colback=teal!4,colframe=teal!35!black,title={Before you move on}]
What does \zh{大衣} mean? (\textquotedblleft{}Overcoat; coat\textquotedblright{}.) Say it once more.
\end{tcolorbox}
\section[\texorpdfstring{\zh{衣} (yī)}{衣 (yī)}]{\zh{衣} — a dot and a bar, then the body}

\label{lesson:ZH-W165-yi-clothes}



\begin{quote}
The word first, then the character inside it.

\begin{itemize}
\raggedright
  \item \emph{Recall:} say \emph{dàyī} — \textbf{R1}, one lesson back, before it decays
\end{itemize}

\begin{quote}
\zh{大} \zh{上}
\end{quote}

\emph{Read it:} each one aloud, with nothing said first — \emph{dà}, \emph{shàng}
\end{quote}

\subsection*{Script — a dot and a bar, then the body}
\hlblockfigure{\detokenize{figures/ZH-W165-yi-clothes-filmstrip.pdf}}{How \zh{衣} is written, stroke by stroke}

\begin{quote}
\zh{衣}
\end{quote}

\textbf{One piece, six strokes.}

\textbf{The top, two strokes.} A dot. A long horizontal under it.

\textbf{The body, four strokes.} A long stroke falling to the lower left. A short vertical inside it that flicks up to the right at its foot. A short stroke on the right falling to the left. Last, a long stroke falling to the lower right.

2 + 4 = six.

\subsection*{Your turn}
\begin{itemize}
\raggedright
  \item \emph{Write it:} \zh{衣}, all six, with the model in front of you
  \item \emph{Cover:} write it again from memory, and count the pieces, not the strokes
\end{itemize}

\begin{tcolorbox}[breakable,colback=teal!4,colframe=teal!35!black,title={Before you move on}]
How many strokes in \zh{衣}? (\textbf{Six}.) What comes first? (\textbf{A dot, then a horizontal}.) Which stroke flicks up? (\textbf{The short vertical inside the body}.)

Source: \href{https://raw.githubusercontent.com/chanind/hanzi-writer-data/68d10a4b21150cae5e1ebbd223eed289cf32d90c/data/\%E8\%A1\%A3.json}{Hanzi Writer Data \zh{衣}, pinned stroke paths and medians}.
\end{tcolorbox}
\section[máoyī]{máoyī — sweater; jumper}

\label{lesson:ZH-C165-maoyi}



\begin{quote}
Before the new one: say the Mandarin for to have a shave, then the Mandarin for an overcoat.
\end{quote}

\subsection*{You'll want to know: máoyī}
\emph{máoyī} — \textquotedblleft{}sweater; jumper\textquotedblright{}.

\begin{quote}
\emph{wǒ de máoyī hěn hǎokàn} — My sweater looks nice.
\end{quote}

Its characters come later; for now it is a word for the ear.

\begin{cousinweb}
\emph{máoyī} is built from \emph{máo} 'hair, wool' and \zh{衣} 'garment', both inherited from Middle Chinese, \emph{maw} and \emph{'j+j} (Guangyun, in Baxter's transcription), which glosses the first as the hair of brows and of beasts.
\end{cousinweb}

\begin{grammarlens}[title={What you've built}]
One more word for everyday things.
\end{grammarlens}

\subsection*{Your turn}
\begin{itemize}
\raggedright
  \item \emph{Say it:} \emph{máoyī}
  \item \emph{Say it:} \emph{máoyī}, once more
  \item \emph{Say it:} \emph{wǒ de máoyī hěn hǎokàn} — My sweater looks nice.
  \item \emph{Recall:} say the Mandarin for to have a shave, then the Mandarin for an overcoat, then say \emph{máoyī} again
\end{itemize}

\begin{tcolorbox}[breakable,colback=teal!4,colframe=teal!35!black,title={Before you move on}]
What does máoyī mean? (\textquotedblleft{}Sweater; jumper\textquotedblright{}.) Say it once more.
\end{tcolorbox}
\section[tiáo]{tiáo — measure word for long things (skirts, scarves, trousers, fish, roads)}

\label{lesson:ZH-C165-tiao}



\begin{quote}
Before the new one: say the Mandarin for an overcoat, then the Mandarin for a sweater.
\end{quote}

\subsection*{You'll want to know: tiáo}
\emph{tiáo} — \textquotedblleft{}measure word for long things (skirts, scarves, trousers, fish, roads)\textquotedblright{}.

\begin{quote}
\emph{wǒ mǎi le liǎng tiáo kùzi} — I bought two pairs of trousers.
\end{quote}

Its characters come later; for now it is a word for the ear.

\begin{cousinweb}
\emph{tiáo} is inherited from Middle Chinese \emph{dew} (Guangyun, in Baxter's transcription), glossed 'a small branch': a long, thin thing. Its use for counting is not given there.
\end{cousinweb}

\begin{grammarlens}[title={What you've built}]
One more word for everyday things.
\end{grammarlens}

\subsection*{Your turn}
\begin{itemize}
\raggedright
  \item \emph{Say it:} \emph{tiáo}
  \item \emph{Say it:} \emph{tiáo}, once more
  \item \emph{Say it:} \emph{wǒ mǎi le liǎng tiáo kùzi} — I bought two pairs of trousers.
  \item \emph{Recall:} say the Mandarin for an overcoat, then the Mandarin for a sweater, then say \emph{tiáo} again
\end{itemize}

\begin{tcolorbox}[breakable,colback=teal!4,colframe=teal!35!black,title={Before you move on}]
What does tiáo mean? (\textquotedblleft{}Measure word for long things (skirts, scarves, trousers, fish, roads)\textquotedblright{}.) Say it once more.
\end{tcolorbox}
\section[qúnzi]{qúnzi — skirt}

\label{lesson:ZH-C165-qunzi}



\begin{quote}
Before the new one: say the Mandarin for a sweater, then the Mandarin for measure word for long things.
\end{quote}

\subsection*{You'll want to know: qúnzi}
\emph{qúnzi} — \textquotedblleft{}skirt\textquotedblright{}.

\begin{quote}
\emph{zhè tiáo qúnzi hěn hǎokàn} — This skirt looks nice.
\end{quote}

Its characters come later; for now it is a word for the ear.

\begin{cousinweb}
\emph{qúnzi} is built from \emph{qún} 'skirt' and the noun ending \zh{子}, both inherited from Middle Chinese, \emph{gjun} and \emph{tsiX} (Guangyun, in Baxter's transcription), which writes the first another way and glosses it as the garment of the lower body. In the word the second loses its tone.
\end{cousinweb}

\begin{grammarlens}[title={What you've built}]
One more word for everyday things.
\end{grammarlens}

\subsection*{Your turn}
\begin{itemize}
\raggedright
  \item \emph{Say it:} \emph{qúnzi}
  \item \emph{Say it:} \emph{qúnzi}, once more
  \item \emph{Say it:} \emph{zhè tiáo qúnzi hěn hǎokàn} — This skirt looks nice.
  \item \emph{Recall:} say the Mandarin for a sweater, then the Mandarin for measure word for long things, then say \emph{qúnzi} again
\end{itemize}

\begin{tcolorbox}[breakable,colback=teal!4,colframe=teal!35!black,title={Before you move on}]
What does qúnzi mean? (\textquotedblleft{}Skirt\textquotedblright{}.) Say it once more.
\end{tcolorbox}
\section[wéijīn]{wéijīn — scarf}

\label{lesson:ZH-C165-weijin}



\begin{quote}
Before the new one: say the Mandarin for measure word for long things, then the Mandarin for a skirt.
\end{quote}

\subsection*{You'll want to know: wéijīn}
\emph{wéijīn} — \textquotedblleft{}scarf\textquotedblright{}.

\begin{quote}
\emph{wàibian hěn lěng nǐ yào dài yì tiáo wéijīn} — It's cold outside, so take a scarf.
\end{quote}

Its characters come later; for now it is a word for the ear.

\begin{cousinweb}
\emph{wéijīn} is built from \emph{wéi} 'to go round' and \emph{jīn} 'a cloth'. The second is inherited from Middle Chinese \emph{kin} (Guangyun, in Baxter's transcription); the first is listed there as \emph{hjw+j}, glossed 'to surround': a cloth that goes round the neck.
\end{cousinweb}

\begin{grammarlens}[title={What you've built}]
One more word for everyday things.
\end{grammarlens}

\subsection*{Your turn}
\begin{itemize}
\raggedright
  \item \emph{Say it:} \emph{wéijīn}
  \item \emph{Say it:} \emph{wéijīn}, once more
  \item \emph{Say it:} \emph{wàibian hěn lěng nǐ yào dài yì tiáo wéijīn} — It's cold outside, so take a scarf.
  \item \emph{Recall:} say the Mandarin for measure word for long things, then the Mandarin for a skirt, then say \emph{wéijīn} again
\end{itemize}

\begin{tcolorbox}[breakable,colback=teal!4,colframe=teal!35!black,title={Before you move on}]
What does wéijīn mean? (\textquotedblleft{}Scarf\textquotedblright{}.) Say it once more.
\end{tcolorbox}
